\subsection{A two-point constant image at arbitrary state count}

\begin{theorem}\label{C:two-constant-image-bisets}
Let $S$ be a finite nonempty $(N,M)$-biset.  Suppose the right
$M$-action has exactly the three transformations
\[
 \{\id_S,c_0,c_1\},
\]
where $c_i$ is constant at a distinct point $i\in S$.
Then its two-object collage cannot have type $\mathbb CP^2$.
The same holds with the sides interchanged.  There is no bound
on the numbers of states or corner elements.
\end{theorem}
\begin{proof}
Commutation with $c_0,c_1$ forces every $N$-actor to fix both
points.  Put $U=\{0,1\}$ and write $X_U\subseteq X_S$ for
the two collage spaces.  Let $\chi:M\to\{0,1\}$ distinguish
identity and constant transformations.  As a right module,
the augmentation module $L=\ker(\mathbb QS\to\mathbb Q)$
is a direct sum of copies of $\mathbb Q_\chi$.
The vector $u=[1]-[0]$ is fixed by $N$.  With $V=L/\mathbb Qu$
as a left $\mathbb QN$-module, there is an extension
\begin{equation}\label{C:two-constant-image-extension}
 0\longrightarrow\mathbb Q\xrightarrow{u}L
 \longrightarrow V\longrightarrow0
\end{equation}
with class $\gamma\in\operatorname{Ext}^1_{\mathbb QN}(V,\mathbb Q)$.
Here $L$ in this display denotes its left-module structure;
its commuting right-module structure is scalar $\chi$.

The coefficient extension
$0\to\mathbb QU\to\mathbb QS\to V\otimes\mathbb Q_\chi\to0$
is the pushout of \eqref{C:two-constant-image-extension}
tensored with $\mathbb Q_\chi$ along the actual inclusion
$u:\mathbb Q_\chi\to\mathbb QU$.
Thus the reduced relative chains of $(X_S,X_U)$ are
\[
 \left(
  \mathbb Q\otimes^{\mathbf L}_{\mathbb QN}V
  \otimes_{\mathbb Q}
  \mathbb Q_\chi\otimes^{\mathbf L}_{\mathbb QM}\mathbb Q
 \right)[1].
\]
The degree-zero homology of the right factor is zero because
$\chi$ takes the value zero.  The K\"unneth summand
\[
 \operatorname{Tor}^{\mathbb QN}_1(\mathbb Q,V)
 \otimes
 \operatorname{Tor}^{\mathbb QM}_1(\mathbb Q_\chi,\mathbb Q)
 \ \subseteq\ H_3(X_S,X_U;\mathbb Q)
\]
therefore has the same boundary computation as in
Lemma~\ref{C:three-state-residual-H2zero}.
More explicitly, the augmentation extension of $\mathbb QU$
identifies the ordinary core act-cone complex with
$(\mathbb Q_\chi\otimes^{\mathbf L}_{\mathbb QM}\mathbb Q)[1]$.
Collapsing the trivially acting $N$-corner induces an actual
$H_2$ isomorphism on $X_U$ by
Lemma~\ref{C:trivial-biset-side-H2}.  Under this identification,
the boundary of the displayed summand is
\[
 \pm\gamma\otimes\id:
 \operatorname{Tor}^{\mathbb QN}_1(\mathbb Q,V)
 \otimes H_2(X_U;\mathbb Q)
 \longrightarrow H_2(X_U;\mathbb Q).
\]
This is a chain-level connecting map: in the splitting
$\mathbb QS=\mathbb QU\oplus V$ the only off-diagonal action
faces come from the extension \eqref{C:two-constant-image-extension};
the right faces are scalar $\chi$, and both vertex augmentations
vanish on the chosen augmentation-module lifts of $V$.

If $\gamma\ne0$, the duality
\[
 \operatorname{Ext}^1_{\mathbb QN}(V,\mathbb Q)
  \cong
 \operatorname{Hom}_{\mathbb Q}
 \bigl(\operatorname{Tor}^{\mathbb QN}_1(\mathbb Q,V),\mathbb Q\bigr)
\]
makes this boundary onto.  Hence $H_2(X_U;\mathbb Q)\to
H_2(X_S;\mathbb Q)$ is zero.  Every degree-two cohomology
class of $X_S$ is relative to $X_U$; the quotient is the
suspension of the corresponding edge-bar pair.  Its positive
products vanish, excluding a projective plane.

If $\gamma=0$, choose an $N$-linear retraction
$q:L\to\mathbb Q$ with $q(u)=1$, and set
$q(s)=q([s]-[0])$.  Then $q(0)=0$, $q(1)=1$, and $q(ns)=q(s)$.
Choose the single rounding function
$\rho:\mathbb Q\to\{0,1\}$ with $\rho(1)=1$ and
$\rho(x)=0$ for $x\ne1$.  The map $r(s)=\rho(q(s))$ fixes
$U$, is $N$-equivariant, and commutes with every right
$M$-actor, since those actors are identity or constant at
$0,1$.  It is an actual biset retraction.
If $|S|>2$, Lemma~\ref{C:proper-biset-retract-CP2} would make
$X_U$ a projective plane whenever $X_S$ were one, contrary
to Theorem~\ref{C:two-cross-arrows-universal}.  The case
$|S|=2$ follows from that theorem directly.
\end{proof}

\begin{remark}\label{C:two-constant-image-scope}
The nonsplit case in the proof gives zero rational
degree-two products without a rank or finiteness assumption.
The split case gives the stated projective-plane exclusion
by an actual retract; no general zero-product assertion is
needed there.
\end{remark}
